\section{Operador bicapa y jerarquía infinita}
\label{sec:jerarquia-espectral-infinita-rev7}
\index{torres!jerarquía infinita}\index{semigrupo espectral}

Sean
\[
 x=A_{\rm rad},\qquad y=C^*_{\rm rad},\qquad
 R_{\rm ch}^2=I,\qquad
 P_\pm^{\rm ch}=\frac{I\pm R_{\rm ch}}2.
\]
El lector de las dos hojas se reúne en el operador
\[
 T_{\rm ang}=\exp(-xI-yR_{\rm ch})
 =q_+P_+^{\rm ch}+q_-P_-^{\rm ch},
\]
con
\[
 q_-=e^{-x+y},\qquad q_+=e^{-x-y}.
\]
Sus momentos espectrales par e impar son
\[
 c_n=\operatorname{Tr}(T_{\rm ang}^n)=q_-^n+q_+^n,\qquad
 s_n=-\operatorname{Tr}(R_{\rm ch}T_{\rm ang}^n)=q_-^n-q_+^n.
\]
La involución \(C^*\mapsto-C^*\) conserva \(c_n\) e invierte \(s_n\).
Equivalentemente,
\[
 c_n=2e^{-nx}\cosh(ny),\qquad
 s_n=2e^{-nx}\sinh(ny).
\]
Ésta es la convención fijada en el capítulo de canales. Para declarar sin
homonimia la presentación alternativa de las fuentes de torres, ponemos
\[
 R_{\rm tow}:=-R_{\rm ch},\qquad
 P_\pm^{\rm tow}=P_\mp^{\rm ch},
\]
\[
 T_{\rm tow}:=e^{-xI+yR_{\rm tow}}=T_{\rm ang},\qquad
 \operatorname{Tr}(R_{\rm tow}T_{\rm tow}^n)
 =-\operatorname{Tr}(R_{\rm ch}T_{\rm ang}^n)=s_n.
\]
Las dos escrituras describen el mismo lector escalar y no constituyen una
segunda orientación de las torres.

\begin{theorem}[Adición, norma y recurrencia]
\label{thm:recurrencia-bicapa-rev7}
Para \(m,n\ge0\),
\[
 c_{m+n}=\frac{c_mc_n+s_ms_n}{2},\qquad
 s_{m+n}=\frac{s_mc_n+c_ms_n}{2},
\]
\[
 c_n^2-s_n^2=4e^{-2nx}.
\]
Si \(p=q_-q_+=e^{-2x}\), ambas sucesiones satisfacen
\[
 X_{n+2}=c_1X_{n+1}-pX_n.
\]
Además,
\[
 \partial_Ac_n=-\frac{n\pi}{180}c_n,
 \qquad
 \partial_As_n=-\frac{n\pi}{180}s_n,
\]
\[
 \partial_{C^*}c_n=\frac{n\pi}{180}s_n,
 \qquad
 \partial_{C^*}s_n=\frac{n\pi}{180}c_n.
\]
\end{theorem}

\begin{proof}
Las identidades de adición resultan al expandir
\((q_-^m\pm q_+^m)(q_-^n\pm q_+^n)\). La norma es
\[
 (q_-^n+q_+^n)^2-(q_-^n-q_+^n)^2
 =4(q_-q_+)^n.
\]
La última igualdad es la recurrencia determinada por las raíces
características \(q_-\) y \(q_+\). Las ecuaciones diferenciales se obtienen
al derivar las expresiones hiperbólicas; el factor \(\pi/180\) pertenece a
la carta angular y no a una selección de especie.
\end{proof}

\begin{theorem}[Semigrupo completo de alturas]
\label{thm:semigrupo-alturas-rev7}
Las alturas de la jerarquía transversal forman
\[
 \mathfrak S
 =\langle90,120\rangle
 =\{90a+120b:a,b\in\mathbb Z_{\ge0},\ a+b>0\}.
\]
Equivalentemente,
\[
 \boxed{\mathfrak S=\{90,120\}\cup\{30r:r\ge6\}.}
\]
\end{theorem}

\begin{proof}
Dividiendo por \(30\), se obtiene el semigrupo
\(\langle3,4\rangle\). Sus elementos positivos son
\(\{3,4\}\cup\{r:r\ge6\}\): \(5\) es la única laguna positiva posterior a
los generadores, y todo \(r\ge6\) se obtiene por inducción sumando \(3\) o
\(4\). Multiplicar de nuevo por \(30\) da la fórmula.
\end{proof}

La presentación conserva también la genealogía de una altura:
\[
 \mathfrak S\cong
 \bigl(\mathbb Z_{\ge0}^2\setminus\{(0,0)\}\bigr)
 \big/\bigl((a+4,b)\sim(a,b+3)\bigr),
\]
pues \(4\cdot90=3\cdot120=360\). La altura no determina por sí sola la
composición que la produjo; las genealogías \(4\cdot90\) y \(3\cdot120\)
permanecen distintas en el registro cronológico enriquecido aunque evalúen
el mismo índice.

Para hacer explícita la recurrencia muestreada en el paso común de treinta,
definamos
\[
 a_{30}:=q_-^{30},\qquad b_{30}:=q_+^{30},\qquad
 u_r:=s_{30r}=a_{30}^r-b_{30}^r.
\]
Entonces
\[
 \boxed{u_{r+2}=(a_{30}+b_{30})u_{r+1}-a_{30}b_{30}u_r.}
\]
El término auxiliar \(u_5=s_{150}\) puede intervenir al propagar la
recurrencia, pero \(150\notin\mathfrak S\); una identidad de cálculo no crea
una nueva altura admisible.

En particular,
\[
90,120,180,210,240,270,300,330,360,390,420,450,480,\ldots
\]
son las primeras alturas. Las ocho alturas destacadas en exposiciones
anteriores son primeras lecturas de \(\mathfrak S\), no su definición ni
su límite.

Como guía humana, su genealogía inicial puede leerse sin convertirla en el
dominio completo:
\[
\begin{array}{c@{\quad}c@{\quad}l}
90&3\cdot30&\text{generador mínimo}\\
120&4\cdot30&\text{generador mínimo}\\
180&2\cdot90&\text{primer modo par}\\
210&90+120&\text{modo compuesto}\\
240&2\cdot120&\text{segundo modo par}\\
270&3\cdot90&\text{repetición de la rama 90}\\
360&3\cdot120&\text{repetición de la rama 120}\\
420&2(90+120)&\text{retorno compuesto}
\end{array}
\]
La última columna registra procedencia estructural; no introduce parámetros
ni afirma que esas ocho filas agoten la jerarquía. Los mismos enteros
\(90,120\) son cardinalidades de incidencia en \(\mathcal F_6\) y
\(\mathcal F_8\), pero esa coincidencia no identifica el operador bicapa con
la inclusión \(\iota_{6,8}\).

La interfaz de los dos generadores es sólo el par
\[
 \mathcal K_{\rm ang}^{90,120}(q_+,q_-):=(s_{90},s_{120}).
\]
No debe confundirse con las resolventes de dos subcolas ni con el funcional
posterior \(K_{6\to8}=12\Delta S_{90}-\Delta S_{120}\), cuya definición y
prueba residen en el \cref{mdg:chap:barbero}. Allí las subcolas viven en la
completación analítica de la jerarquía; no son nuevos generadores ni
coeficientes de especie.

\begin{definition}[Módulo convergente de refinamientos espectrales]
\label{def:modulo-refinamientos-torre-rev7}
El dominio integral de la elevación global es
\[
 \mathcal E_{\mathfrak S}
 :=\left\{\eta\in\mathbb Z^{\mathfrak S}:
 \sum_{h\in\mathfrak S}|\eta_hs_h|<\infty\right\},
 \qquad s_h=q_-^h-q_+^h.
\]
Cada \(\eta\in\mathcal E_{\mathfrak S}\) determina el carácter espectral
\[
 \mathcal R_\eta
 :=\exp\!\left(\sum_{h\in\mathfrak S}\eta_hs_h\right).
\]
\end{definition}

La ley de masas no reside en esta definición: su firma central
\((n_A,s_C,k_{\Delta_4},b_{120},b_\zeta)\in\mathbb Z^5\) se construye antes
desde el registro. La elevación por \(\eta\) actúa después y conserva la
procedencia del término básico \(b_{120}s_{120}\). En particular,
\(N_{120}=b_{120}+\eta_{120}\) es el coeficiente total evaluado, no una
nueva coordenada de firma ni una descomposición inversa única.

La completitud del semigrupo no selecciona por sí sola una especie. Cuando
\(\eta\) se obtiene desde un residual metrológico suministrado, se tiene una
reconstrucción calibrada exacta; la fuerza prospectiva requiere que la ruta,
el acarreo y el devanado produzcan \(\eta\) antes de abrir ese residual. La
recurrencia anterior genera todos los momentos una vez fijado el
refinamiento, sin introducir una ley distinta en cada altura.
